\chapter[Sermon 78. Counsel Is Given to the Godly, Which Are Commanded to Go Out of Babylon. Enemies Are Stirred Up Against Babylon, and They Are Commanded Not to Spare Her.]{Counsel Is Given to the Godly, Which Are Commanded to Go Out of Babylon. Enemies Are Stirred Up Against Babylon, and They Are Commanded Not to Spare Her.}
\label{sermon:078}
\markright{Sermon 78. Counsel Is Given to the Godly, Which Are Commanded to ...}

And I heard another voice from Heaven saying, “Come away from her, my people, that you be not partakers of her sins, lest you receive of her plagues. For her sins have gone up to Heaven, and the Lord has remembered her wickedness. Give her reward even as she rewarded you, and give her double according to her works. And pour into her double in the same cup which she filled unto you. And as much as she glorified herself and lived wantonly, so much power she will have in punishment, and sorrow, for she says in her heart, ‘I sit being a queen, and am no widow, and shall see no sorrow.’ Therefore her plagues shall come in one day, death and sorrow, and hunger, and she shall be burned with fire, for strong is the Lord God who shall judge her.”

\index{Papacy}\index{Rome}The second point of this chapter is the faithful counsel of the Lord, given to the godly, concerning how they should behave during the prosperity and destruction of the city. Rome has indeed for a long time been mistress of the world; the riches and pleasures of the entire world have been seen at Rome. If anyone at Rome or in the provinces showed himself compliant and obedient to the Romans, and loved the Roman religion greatly, and conformed to the corrupt manners of the Romans, he was highly esteemed and could, as it were, by degrees, attain to high promotion and dignity, to the greatest riches, and the most desirable pleasures. If any man resisted the Roman religion and would not assent to the Romans, he was afflicted with persecution, plundered and driven into exile, imprisoned, or led to execution. Therefore, the godly were grievously tempted and did not know whither to turn. As we see the like done today in New Rome and the papist kingdom throughout the world. Wherefore, God, who does not will that man should perish but be saved, gives here the best counsel of true felicity and salvation; those who obey it are blessed.

And directly he shows the author of this counsel from the beginning, so that he might establish its authority, and that we might boldly receive it. “I heard,” he says, “another voice from heaven; therefore this counsel proceeds from God.” Those who follow it obey God; those who do not obey it despise and reject God’s counsel. And what is this counsel? Briefly, it is plain, possible, honest, and wholesome, without doubt. “Come away,” says the Lord, “from her—that is, Babylon, both old and new—my people: you who will be called the people of God and be written in the number of the citizens of God.” This is his counsel, and no other. God gave this same counsel to his ancient people through his prophets when they were in captivity in Babylon. For Isaiah says in chapters 48 and 52, “Depart, depart, come away from there; do not touch anything unclean.” “Come away from her, be cleansed, you who bear the vessels of the Lord.” And Jeremiah in chapter 51, “Flee from the midst of Babylon, and let every man save his soul, so that he will not be rooted out in her wickedness.” For the time of God’s vengeance is at hand; he will reward her. Therefore, the Lord counsels us to flee, so that our souls may be saved. For otherwise, unless we flee, we shall perish. However, the prophets did not teach the Israelites to flee from Babylon physically, by local motion, as they term it. For Jeremiah in chapter 29 exhorts the captive people to dwell in Babylon and to make their provision there until the time of deliverance arrives. For then they must come out of Babylon. In the meantime, he would have them depart not by physical motion, but by difference of manners. For although they shall dwell in the midst of the superstitious, ungodly, and idolaters, the Lord does not want them to be like them. Therefore, this fleeing consists in abstaining and refraining from ungodliness, idolatry, sins—namely, bloodshed, usury, pride, lechery, and other like vices—but persevering in true godliness and innocence.

In like manner, if the godly were to flee under the old Roman Empire, they would everywhere fall again into the hands of the Romans; as we also at this day, although we change our place, have papistry either ever-present or imminent. Therefore, the Apostle speaks well: we must get ourselves out of the world if we are to be conversant with sinners. This, therefore, is the true and godly flight: if remaining in this world bodily, in mind and manners we depart furthest out of the world, so that we abstain from all idolatry and profane worship, if we do not allow it, if it does not please us; if we neither assent nor conform ourselves to the manners of the ungodly; if we shall not betray our religion, either for people or for worldly gain. So therefore, the Christians who lived under the Roman Empire fled Rome, utterly abstaining from the worship of idols and the corrupt manners of the gentiles, although they lived among the heathens. For that the ancient churches in Asia were such, we have heard in the second and third chapters of this book. Albeit therefore that we also dwell under the Popish kingdom and empire that persecutes the gospel, yet must we flee papistry, that is to say, Popish churches; none of the godly ought, for worship or obedience’ sake, to enter in, to acknowledge, allow, or use any Popish rites or ceremonies, but from their vices and corruptions to flee as far as is possible. For so the Apostolic scripture teaches us in Romans 12, 2 Corinthians 6, Ephesians 5, and 1 Peter 4. And John, as it were, expounding himself, says, “Be not partakers of her sins, lest you share in her plagues.” And sins are not only those which are done against the second table, but also those that are committed, and that much more against the first table; of which sort are idolatry, impiety, the abuse of God’s holy name, strange worship, against the second and fourth precepts of the first table. Those were then, and so are at this day taken for very good works, where they are abominations. Partaking is chiefly in the communion of sacred things, again if they are given to the same dissolute riot with filthy men. If therefore we beware of those things, we flee out of Babylon and follow the good counsel of God.

But in this we commonly offend today, those who are called Gospel advocates. For many think it sufficient if they observe some religion in their heart privately, and openly associate with them, which may either help or hurt. You will find those who kneel before idols, who hear Mass and popish services. There are some who know many abominations of the popish priesthood, but nevertheless make their sons priests, because the promotions, and the clerical life—that is to say, the wealthy and pleasant life—appeals to them. There are some who, through the bond of marriage, force their children into the midst of Papistry; and these regard nothing else but riches, worldly honors, and friendships. Against all these, the prophets and the Apostles, and at this present Christ the Son of God from the right hand of the Father, thunder and cry aloud, “Come away from her, my people, and be not partaker with her sins.” These words do not admit of any witty or civil reasoning, nor carnal or crafty qualifying. For it follows, lest they receive of her plagues. For if the like Rome, if the like the Roman religion, if the Roman prelacy please you, riches and promotions, if the Roman corruption content you, let the judgment, pain, and damnation due to Antichristianism content you also.

We also have, at this time, what answer we may give to Romanists, objecting and claiming that we are rebelling or committing apostasy, and for the same reason, the crime of schism. “You have fallen from the holy church of Rome,” they say, “and by that very abandonment, you declare openly that you are sectaries and schismatics.” To this we answer that we make a distinction within the church of Rome. For we acknowledge an old church of Rome, venerable and apostolic. Of which Paul wrote: “Your faith is proclaimed throughout the world.” Whoever departs from that church, without doubt, will be both a schismatic and will perish forever.

\index{Repentance}There is again another church of Rome, new, and entirely contrary to the old, no longer apostolic but rather Papistic, wherein are not ministers of the word and sacraments, but either princes, not unlike the Gentiles, or merchants, from whom the sacraments, the remission of sins, heaven itself, and all things in the church are sold for a little money. They teach a doctrine that deviates entirely from the doctrine of the gospel. These are openly, not infected, but swimming and stinking of the most shameful vices, even the filthiness of whoredom—to say nothing now of Christian bloodshed. And there is no sign of repentance in them. To persevere with these, to commune with these, is to perish eternally. Therefore, from the company of these men, the Lord commands us here to depart, yes, and to flee from. That is what we have done, at the Lord’s commandment, which openly here commands us to come away, depart, and flee from the scarlet woman and from this Babylon.

There are also other notable passages commanding this departure, which whoever wishes to know and consider should read Deuteronomy 13, Jeremiah 23, the words of the Lord in the gospel of Luke 6:23 and 24, and Matthew 7:23. Read also both epistles of Paul to Timothy, especially 6:3-5 of the first, and 1:3 and 4, and 2:4 of the second. In Romans 16, he says: “I beseech you, brethren, mark those who cause divisions and put up to evil, contrary to the doctrine which you have learned, and avoid them. For those who are such do not serve our Lord Jesus Christ, but their own appetites, and through flattering words.”

And explaining the reason why we should flee from Babylon, he speaks of profit and loss. Lest we receive her plagues. For whoever aligns himself with the ungodly, idolaters, and those who are filthy and unclean, receives the same reward as they do: and the reward of this present life is a curse, a reprobate mind, and various calamities, as described in chapter 16 and elsewhere, and after this life, everlasting torments. Therefore, he is not speaking of a trivial matter when he warns of fleeing from Babylon or avoiding the Roman religion. Many believe these things, for they do not consider how great is the abomination of the church of Rome before God; and therefore they hear these things as a fable and persevere in the same kind of life in which they are and have lived until now. But he does not lie who says that those who do not prepare to flee from Babylon will shortly perish with Babylon and with the entire fellowship of the wicked. Woe to them.

However, since the wicked in this world are commonly fortunate (and many from this infer that God does not know our affairs, or at least if he knows them, does not care greatly), the Apostle, or an oracle brought from heaven, adds: “For her sins have risen up to heaven, and the Lord has remembered her wickedness.” God truly never forgets iniquities. For all things are evermore present before him. Yet he seems not to remember when he does not punish. For so men suppose: but when he punishes and visits sinners, he seems utterly to have had consideration of our affairs, and to have remembered wickedness and wicked men. Therefore God is righteous, and mindful of evil, and of good also: and when he sees that the time will recompense all men's works, and chiefly the evil. In the meantime, he signifies also that the sins of old and new Rome are great and full of enormity. For in Genesis 19, the sins of Sodom are said to have risen up to heaven, as if to cry out against those who committed them and demand vengeance. So we read in Jeremiah 51 that the sins of Babylon ascended up to the clouds. For John, in a manner everywhere, uses the places of Scripture, to the intent he might give his book more authority, although otherwise inspired by the Holy Spirit. Indeed, the old satirical poets, such as Horace, Juvenal, and others, wrote severely against the sins and vices of old Rome. There remain also at this day many sharp writings against Rome, and the Cardinals and Prelates of the Roman church, and Pasquils innumerable (Pasquil at this day is a satirical writer, one in place of many), so that as well at this day as in times past the sins of Rome cry up to heaven itself.

He proceeds after this to recount again the plagues and most certain destruction of Rome, which is the third section of this chapter. Here, the most horrible and cruel manner of destruction and subversion is exceedingly well described. For God is portrayed as calling on and exhorting the soldiers, the commissaries, and the executors of his judgment to vengeance: and that they should punish her most extremely, and spare her not, but reward her most abundantly, and mete unto her by the same measure wherewith Rome has measured to others. For here takes place that saying of the Lord, a common saying with all nations: with the same measure wherewith you mete, others shall mete unto you again, and there shall be given good measure, pressed, shaken, and running over. Therefore, saying that Rome has robbed the whole world and seduced the whole world, rightly and by the just wrath of God was she spoiled and utterly subverted. The Goths accomplished these things with great faith and diligence, so that we cannot doubt but that New Rome, and that Apostate See, must of her enemies, whom the Lord has prepared, and of the Angels gathering the tares, be plucked all to pieces. And what shall become of her in another world, we may gather hereof, that he beats in so often that her evils shall be doubled without mercy, her pain also, mourning, and grievous torments. These things doubtless are grievous and horrible. Would God they might be perceived by faithful minds. And again, this passage is written as it were word for word from the 50th chapter of Jeremiah: where you read to this effect: be avenged on Babylon, and as she did, do the unto her. Spoil and destroy says the Lord, and accomplish all that I have commanded thee. Destroy her, that nothing remain. Entrench round about, that no man escape. Reward her after her work: and according to all things that she hath done, do the unto her. For she has been proud against the Lord, and against the holy one of Israel. Thus said the Lord in Jeremiah. You see therefore where the Lord borrows his own at this present. You see what every city, or commonwealth, or man may promise himself, if being enriched by the loss of others, he live voluptuously and proudly in this world. For God is the same always, and his judgments are keen against all the ungodly.

And he has been involved in the causes of ruin, cruelty, covetousness, extortion, slaughter, burning, with which Rome has laid waste the entire world. But he proceeds more expressly to recite other causes: namely, pride, glorying, and boasting, security, riot, pleasures, and indulgence. For it follows: as much as she has glorified herself and lived luxuriously, and so forth. And again: for in her heart she says, “I sit as a queen,” and so forth. He has borrowed these things also from Isaiah 47. Where Babylon glories thus also, and with so many words. Rome in times past gloried herself to be Lady of the world, and that she should be everlasting. For they stamped on silver coins the words “Rome eternal.” They had thought that the kingdoms would never be taken from her. She thought therefore that she would never be a widow. And I doubt not but the Germans borrowed from the Romans that Germanic word \textit{Römer}, by which they mean to boast or brag stoutly, which seems to have been peculiar and proper to the Romans. She was careless and secure. She had not thought to be overthrown. She said, “I shall see no morning; I will have no morning cheer; I will always sing, ‘Gaudeamus.’” The Romanists at this day also boldly make their boast that no emperors, no kings, no people, no heretics, or schismatics (for so they term the enemies of Roman wickedness, godly and learned men) have yet successfully assailed Rome. That the enemies of the church of Rome have always been oppressed, that she has always triumphed over her enemies, these seven or eight hundred years and more. That the ship of St. Peter may be sorely troubled, tossed, and overwhelmed with waves and billows, but cannot be drowned; and therefore that the See of Rome shall be a perpetual queen and lady of all realms and churches. and so forth.

\index{Arethas of Caesarea}But hear now God’s judgment: for as much as she is proud, vain, glorious, careless, and wicked, her plagues will come in one day. Arethas notes that by “one day” is signified a sudden destruction, and that she will perish when she least expects it. He recounts her plagues in order: death, mourning, famine, and fire. Historical accounts testify that these things were fulfilled in old Rome, as I have spoken of before. Therefore, we doubt nothing at all, but that new Rome also will be torn apart by humans and by God’s angels and uprooted. And lest any man should think this impossible (for great is the power and majesty of Rome, so that anyone who had said in John’s time that Rome would fall would have seemed to speak something as impossible as saying the sky would fall), he adds immediately: for the Lord God who will judge her is mighty. Therefore, let us not doubt the fall of Papistry. For the Lord is true, just, and almighty. To whom be glory forever and ever. Amen.
